\section{An executable witness: constructing what the successor needs}\label{sec:study}

The framework should produce something that works. We therefore accompany the
paper with a finite study that constructs new executable instruments, checks their
outputs, and asks when the complete means of using them are affordable and
available. Its purpose is to expose distinctions that a compelling diagram could
hide: a finished instrument can fail on a successor; a reusable constructor can
build a valid replacement; current rights and qualification remain necessary; and
retention can impose obligations that make the next outcome harder to reach.

\begin{statusbox}
\textbf{What is supplied and what is executed.} The task language, compiler,
specifications, alternative routes, rights scenarios, and planning costs are
authored. The generated programs and their output checks are executed. The
constructor is \emph{supplied, not learned}. No physical sensor, customer,
transaction, model call, or autonomous discovery is represented.
\end{statusbox}

\subsection{Constructing a fresh instrument}

A task specifies named integer input fields $f_j$, affine calibration coefficients
$a_j,b_j$, aggregation weights $w_j$, and a threshold $\theta$. Its required
measurement and decision are
\begin{equation}
 z(x)=\sum_{j=1}^{d}w_j\bigl(a_jx_{f_j}+b_j\bigr),
 \qquad d(x)=\mathbf 1\{z(x)\geq\theta\}.
 \label{eq:study-instrument}
\end{equation}
The compiler combines these terms into a standalone Python program. It consumes
JSON records, requires the exact integer schema, computes the measurement and
Boolean decision, and emits a specification hash. This is actual source
generation from a supplied specification, using ordinary algebra. It is not a
discovery of a new measurement law.

The source instrument computes $2a+2b+c+1$. The first successor instead requires
$2q-p+6r+1$, with different field names, coefficients, and threshold. Further
successors require four and five fields. The finished source program cannot
consume these schemas. Keeping that program alone therefore does not solve the
successors. The retained generic compiler creates three fresh programs from their
specifications; their correctness does not depend on executing the finished donor
program.

For each task we enumerate every vector in
$\{-2,-1,0,1,2\}^{d}$. A separate calculation checks both the numerical result and
the decision. A strong ordinary control directly interprets the same declarative
specification, without needing to generate a program. Table~\ref{tab:study-programs}
reports all cases. The output field named \texttt{accepted} denotes the
per-record threshold decision; commercial acceptance is not implied. Program
qualification means matching the specified result on the entire declared test
domain, including correctly negative decisions.

\begin{table}[htbp]
\centering\small
\begin{tabularx}{\linewidth}{YrrP{1.58in}}
\toprule
Task & Input cases & Both methods correct & Finished source program\\
\midrule
Source: three fields & 125 & 125 & Same task\\
Successor 1: three fields & 125 & 125 & Schema rejection\\
Successor 2: four fields & 625 & 625 & Schema rejection\\
Successor 3: five fields & 3,125 & 3,125 & Schema rejection\\
\midrule
Total & \textbf{4,000} & \textbf{4,000} & Three incompatible successors\\
\bottomrule
\end{tabularx}
\caption{Actual generated programs and the ordinary interpreter agree with the
specifications on the complete stated finite domain. These authored fresh tasks
are not independently selected natural tasks. No runtime or learning advantage is
claimed.}
\label{tab:study-programs}
\end{table}

\subsection{Completing the means of delivery}

Correct output is one component of the modeled outcome. An internal delivery
also requires a reusable testing fixture, current access terms, fresh
qualification, and production resources. An external complete service is a
separately available conventional route under its own supplied standard terms.
Missing internal consent does not disable that external alternative.

The fresh internal route costs 54 planning units and takes eight time units:
18 for construction, 10 for the fixture, 6 for current terms, 8 for qualification,
and 12 for production. Inheriting the fixture removes its repeated cost and time,
giving 44 units and six time units. Inheriting the supplied constructor additionally
changes assumed fresh construction cost from 18 to 7 and its duration from three
to one, giving 33 units and four time units. The complete external service costs
58 units and takes four time units. These are \emph{synthetic planning
assumptions}, not measured savings from compilation, labor estimates, currency, or
market quotations.

The constructor state also carries a current maintenance obligation $m$. This
must be paid even if the external service or nonentry is selected. Every delivery
expense precedes its modeled receipt of 80 units. A route is feasible only if
every cash prefix remains nonnegative, the deadline is met, and its current
rights and qualification conditions hold. Constructor acquisition and full
continuation costs are accounted for separately in the economic analysis; the
state comparison here conditions on what has already been retained.

\subsection{The frontier expands---and can contract}

Before execution, the protocol fixed cash values
$\{32,33,43,44,53,54,57,58\}$, deadlines $\{3,4,5,6,7,8\}$, two internal-consent
settings, and maintenance values $\{0,12,26\}$. Their Cartesian product contains
288 cells. Evaluating the three inherited states gives 864 complete-route
decisions. Each decision selects the least-cost feasible supplied route, or
nonentry when none qualifies. This is exact only within the declared catalogue.

\begin{table}[htbp]
\centering\small
\begin{tabularx}{\linewidth}{P{.92in}rYYY}
\toprule
Maintenance $m$ & Cells & Fresh state & Fixture-only & Constructor state\\
\midrule
0 & 96 & 12 & 22 & \textbf{40}\\
12 & 96 & 12 & \textbf{22} & 20\\
26 & 96 & 12 & \textbf{22} & 0\\
\bottomrule
\end{tabularx}
\caption{Number of cells admitting an accepted modeled delivery. The selected
grid is a finite comparison, not a probability distribution. Fresh and fixture-only
counts repeat because the maintenance obligation belongs only to the constructor
state. These unequal-state counts are not a ranking of learning algorithms.}
\label{tab:study-frontier}
\end{table}

The positive witness is transparent. With cash 33, deadline four, current internal
consent, fresh qualification, and zero maintenance, the constructor route delivers.
The fresh, fixture-only, and external routes do not fit those resources. Remove
current rights or qualification and the internal result becomes infeasible.

The adverse witness is equally consequential. With cash 58, deadline four, and
no internal consent, the fixture-only state can buy the external complete service.
Retaining a constructor with maintenance 12 leaves insufficient cash for that
same service. The additional asset has contracted the reachable set. This does
not contradict conditional state dominance: its no-greater-obligations premise
has failed. A useful retained object is not automatically a better productive
state.

\subsection{What the strong control establishes}

An ordinary analytical planner receives the same specifications, constructor
recipe, retained fixture, resource limits, rights, costs, and complete routes. It
agrees with the constructor-state planner in \textbf{all 288 decisions}. Its
ordinary interpreter also matches every generated instrument's required output.
The control is therefore not deprived of the very content responsible for the
conditional gain. Compilation is not technically necessary for these arithmetic
tasks: the interpreter already supplies the required answers. The emitted
programs demonstrate portable construction, not the uniquely necessary or
empirically cheapest implementation.

The study establishes a finite example of \emph{fresh construction using retained
means}. It does not establish that experience improved the constructor, that a
special discovery policy outperformed ordinary engineering, or that the supplied
cost advantage would occur in practice. The conventional tie is a substantive
result: productive value and proprietary methodological superiority are different
claims. The three fresh schemas demonstrate donor-artifact incompatibility; they
do not establish open-ended transfer outside the supplied language.

\subsection{Reproduction and claim boundaries}

From \texttt{docs/learning-to-build/study}, run:
\begin{lstlisting}
python3 study.py --output evidence
python3 check.py --evidence evidence
\end{lstlisting}
The first command generates the instruments, compact results, and all grid cells.
The second imports no producer code: it separately executes the instruments and
replays route costs and cash prefixes. This is a separate-source engineering
cross-check, not an independent institutional evaluation. The retained protocol
and result hashes bind the task definitions and implementation. Both entry
points pin the exact registered protocol; this reproduction is not a general
configurable simulator. The checker recomputes the complete reported record,
including all counts, hashes, flags, and row cardinalities. Normal and
optimized Python execution produced byte-identical reported results.

Adverse checks reject incompatible schemas, a tampered task contract, missing
current qualification, and absent internal rights, while preserving the
maintenance counterexample. Exact checking covers the authored domain and
assumptions. It does not establish real demand, observed consent, physical
qualification, measured economic productivity, or construction learning. The
next experiment must earn those stronger bridges rather than rename this witness.

\begin{thesisbox}
\textbf{The lesson.} Keeping a finished object and keeping a means of building
the next object are different achievements. Either becomes productive only when
the complete current means of accepted use are available.
\end{thesisbox}
